%==============================================================================
% Scale Effects and Jones Fix - Exam Reference
% Endogenous Growth Theory
%==============================================================================

\section{Scale Effects and the Jones (1995) Fix}
%==============================================================================

\subsection*{The Scale Effects Problem}

In the basic Romer (1990) model, knowledge creation exhibits constant returns to the
stock of ideas $M$ and R\&D labor $L_R$:

\[
\dot{M} = \lambda M L_R
\tag{Romer Knowledge}
\]

On a balanced growth path with $L_R \propto L$ (population), this implies:

\[
\gamma = \frac{\dot{M}}{M} = \lambda (1-\alpha)L - \rho
\tag{Romer Growth}
\]

\begin{tcolorbox}[colback=red!10,colframe=red!60,title=Scale Effects Problem]
Growth rate $\gamma$ increases with $L$ (population/labor force).
\bigskip

\textbf{Prediction:} Larger economies should grow faster.
\bigskip

\textbf{Problem (Jones 1995):} Empirically rejected!
\end{tcolorbox}

%-----------------------------------------------------------------------------
\subsection{Jones (1995): Empirical Rejection}

\begin{itemize}
    \item R\&D personnel increased $\sim 10\times$ since 1950s
    \item GDP growth did \textbf{NOT} accelerate
    \item If scale effects held, we should have seen acceleration
\end{itemize}

\begin{tcolorbox}[colback=orange!10,colframe=orange3,title=Key Insight]
The original Romer model generates \textbf{scale effects} because
constant returns to $M$ implies that having more ideas makes it easier
to create new ideas.
\end{tcolorbox}

%-----------------------------------------------------------------------------
\subsection{Jones' Fix: Decreasing Returns to R\&D}

Jones (1995) proposes decreasing returns to the stock of ideas:

\begin{derivationbox}
Original Romer: $\dot{M} = \lambda M L_R$ \quad (CRS to $M$)

Jones Fix: $\dot{M} = \lambda M^\phi L_R$ \quad with $\phi < 1$
\bigskip

Key difference: $\phi < 1$ implies decreasing returns to the stock of ideas.
Having \textbf{more ideas makes it harder} (not easier) to generate new ones.
\end{derivationbox}

%-----------------------------------------------------------------------------
\subsection{Derivation of No-Scale Growth}

\begin{derivationbox}
\textbf{Step 1: Free entry condition}
\[
r = \frac{\partial \pi}{\partial M} = \lambda \alpha M^{\phi-1} L_E
\tag{Arbitrage}
\]

\textbf{Step 2: On BGP, $r = \rho$}
\[
\rho = \lambda \alpha M^{\phi-1} L_E
\tag{BGP Rate}
\]

\textbf{Step 3: Knowledge accumulation}
\[
\frac{\dot{M}}{M} = \lambda M^{\phi-1} L_R
\tag{Knowledge}
\]

\textbf{Step 4: Using BGP condition}
\[
\frac{\dot{M}}{M} = \frac{\rho}{\alpha} \frac{L_R}{L_E}
\tag{}
\]

With $L_R / L_E = (1-\phi) n$ (from resource constraint):
\[
\boxed{\gamma = \frac{n}{1-\phi}}
\tag{Jones Growth}
\]
\end{derivationbox}

\begin{resultbox}
\textbf{Key Result:} $\gamma = \dfrac{n}{1-\phi}$
\bigskip

\textbf{No scale effects!} Growth depends on:
\begin{itemize}
    \item $n$ = population growth rate
    \item $\phi$ = R\&D curvature parameter ($0 < \phi < 1$)
\end{itemize}
\bigskip

\textbf{IMPORTANT:} If $n = 0$, then $\gamma = 0$!
This is called \textbf{semi-endogenous} growth.
\end{resultbox}

%-----------------------------------------------------------------------------
\subsection{Key Implications}

\begin{enumerate}
    \item \textbf{Higher population growth $\Rightarrow$ higher growth}
    \item \textbf{Zero population growth $\Rightarrow$ no growth}
    \item Growth is \textbf{not} endogenous in the same way as Romer (1990)
    \item Scale (size of economy) does \textbf{not} affect long-run growth
\end{enumerate}

\begin{tcolorbox}[colback=blue!10,colframe=blue3,title=Semi-Endogenous Growth]
Jones (1995) model is called "semi-endogenous" because:
\begin{itemize}
    \item Growth rate is determined by exogenous $n$ (population growth)
    \item Level effects from $s$ (savings rate) but no growth effects
    \item Need population growth for sustained growth
\end{itemize}
\end{tcolorbox}

%-----------------------------------------------------------------------------
\subsection{Model Comparison}

\begin{center}
\begin{tabular}{lccc}
\toprule
\textbf{Model} & \textbf{Scale Effects?} & \textbf{Growth Driver} & \textbf{Type} \\
\midrule
Solow & No & Exogenous $g$ & Exogenous \\
AK & No & $s$ (level only) & Endogenous \\
Romer (1990) & Yes ($\gamma \propto L$) & Endogenous ($L$) & Endogenous \\
Jones (1995) & No & $n$ (population growth) & Semi-Endogenous \\
\bottomrule
\end{tabular}
\end{center}

%==============================================================================
\end{document}
